% GENERATED FILE. Edit canonical lessons, then run npm run generate:books.
% canonical-source-hash: fnv1a64:360fe429e9a4cb95
% canonical-lessons: AR-C100-misad, AR-C100-mamarr, AR-C100-maydan, AR-C100-matar, AR-C100-mina

\chapter{Lift, Corridor, Public square, Airport, Port}
\label{ch:ar-lift-corridor-public-square-airport-port}
\hlchaptermodality{100}

\begin{chapteropening}
\textbf{By the end of this chapter:} \emph{I can say a lift, a corridor, a public square, an airport and a port.}

Complete the last lesson of chapter 100: I can say a lift, a corridor, a public square, an airport and a port.
\end{chapteropening}

\section[\texorpdfstring{\ar{مصعد} (miṣʿad)}{مصعد (miṣʿad)}]{\ar{مصعد} (miṣʿad) — a lift}

\label{lesson:AR-C100-misad}



\begin{quote}
Before the new one: say the Arabic for a language, then the Arabic for a story.

\emph{Read it:} \textbf{\ar{ال}}, then say what it means

(\emph{Al-}, \textquotedblleft{}the\textquotedblright{} — the \emph{al-} hiding in \emph{algebra} and \emph{alcohol}.)
\end{quote}

\subsection*{You'll want to know: \ar{مصعد}}
\textbf{\ar{مصعد}} — \emph{miṣʿad} — \textquotedblleft{}a lift\textquotedblright{}.

\begin{cousinweb}
\textbf{\ar{مصعد}} (\emph{miṣʿad}) is built on the native Arabic root \emph{ṣ-ʿ-d}, \textquotedblleft{}to go up, to ascend\textquotedblright{} (Lane). \emph{Miṣʿad} is an instrument noun, \textquotedblleft{}a thing for going up\textquotedblright{}; the lift is a modern sense.
\end{cousinweb}

\begin{grammarlens}[title={What you've built}]
One more word for this chapter.
\end{grammarlens}

\subsection*{Your turn}
\begin{itemize}
\raggedright
  \item \emph{Say it:} \emph{miṣʿad}
  \item \emph{Say it:} \emph{miṣʿad}, once more
  \item \emph{Say it:} say \emph{miṣʿad}
  \item \emph{Recall:} say the Arabic for a language, then the Arabic for a story, then say \emph{miṣʿad} again
\end{itemize}

\begin{tcolorbox}[breakable,colback=teal!4,colframe=teal!35!black,title={Before you move on}]
What does \ar{مصعد} mean? (\textquotedblleft{}A lift\textquotedblright{}.) Say it once more.
\end{tcolorbox}
\section[\texorpdfstring{\ar{ممر} (mamarr)}{ممر (mamarr)}]{\ar{ممر} (mamarr) — a corridor}

\label{lesson:AR-C100-mamarr}



\begin{quote}
Before the new one: say the Arabic for a story, then the Arabic for a lift.
\end{quote}

\subsection*{You'll want to know: \ar{ممر}}
\textbf{\ar{ممر}} — \emph{mamarr} — \textquotedblleft{}a corridor\textquotedblright{}.

\begin{cousinweb}
\textbf{\ar{ممر}} (\emph{mamarr}) is built on the native Arabic root \emph{m-r-r}, \textquotedblleft{}to pass, to go by\textquotedblright{} (Lane). \emph{Mamarr} is the \emph{ma-} place noun, \textquotedblleft{}a place of passing\textquotedblright{}.
\end{cousinweb}

\begin{grammarlens}[title={What you've built}]
One more word for this chapter.
\end{grammarlens}

\subsection*{Your turn}
\begin{itemize}
\raggedright
  \item \emph{Say it:} \emph{mamarr}
  \item \emph{Say it:} \emph{mamarr}, once more
  \item \emph{Say it:} say \emph{mamarr}
  \item \emph{Recall:} say the Arabic for a story, then the Arabic for a lift, then say \emph{mamarr} again
\end{itemize}

\begin{tcolorbox}[breakable,colback=teal!4,colframe=teal!35!black,title={Before you move on}]
What does \ar{ممر} mean? (\textquotedblleft{}A corridor\textquotedblright{}.) Say it once more.
\end{tcolorbox}
\section[\texorpdfstring{\ar{ميدان} (maydān)}{ميدان (maydān)}]{\ar{ميدان} (maydān) — a public square, a field}

\label{lesson:AR-C100-maydan}



\begin{quote}
Before the new one: say the Arabic for a lift, then the Arabic for a corridor.
\end{quote}

\subsection*{You'll want to know: \ar{ميدان}}
\textbf{\ar{ميدان}} — \emph{maydān} — \textquotedblleft{}a public square, a field\textquotedblright{}.

\begin{cousinweb}
\textbf{\ar{ميدان}} (\emph{maydān}) is listed by Lane under \emph{m-y-d} as a race-ground for horses. He quotes two native derivations, from \emph{māda}, \textquotedblleft{}to move\textquotedblright{}, or from \emph{madana}, \textquotedblleft{}to dwell\textquotedblright{}; some modern scholars trace it instead to an Iranian source, so the origin is not settled.
\end{cousinweb}

\begin{grammarlens}[title={What you've built}]
One more word for this chapter.
\end{grammarlens}

\subsection*{Your turn}
\begin{itemize}
\raggedright
  \item \emph{Say it:} \emph{maydān}
  \item \emph{Say it:} \emph{maydān}, once more
  \item \emph{Say it:} say \emph{maydān}
  \item \emph{Recall:} say the Arabic for a lift, then the Arabic for a corridor, then say \emph{maydān} again
\end{itemize}

\begin{tcolorbox}[breakable,colback=teal!4,colframe=teal!35!black,title={Before you move on}]
What does \ar{ميدان} mean? (\textquotedblleft{}A public square, a field\textquotedblright{}.) Say it once more.
\end{tcolorbox}
\section[\texorpdfstring{\ar{مطار} (maṭār)}{مطار (maṭār)}]{\ar{مطار} (maṭār) — an airport}

\label{lesson:AR-C100-matar}



\begin{quote}
Before the new one: say the Arabic for a corridor, then the Arabic for a public square.
\end{quote}

\subsection*{You'll want to know: \ar{مطار}}
\textbf{\ar{مطار}} — \emph{maṭār} — \textquotedblleft{}an airport\textquotedblright{}.

\begin{cousinweb}
\textbf{\ar{مطار}} (\emph{maṭār}) is built on the native Arabic root \emph{ṭ-y-r}, \textquotedblleft{}to fly\textquotedblright{} (Lane). \emph{Maṭār} is the \emph{ma-} place noun, \textquotedblleft{}a place of flying\textquotedblright{}; the airport is a modern sense.
\end{cousinweb}

\begin{grammarlens}[title={What you've built}]
One more word for this chapter.
\end{grammarlens}

\subsection*{Your turn}
\begin{itemize}
\raggedright
  \item \emph{Say it:} \emph{maṭār}
  \item \emph{Say it:} \emph{maṭār}, once more
  \item \emph{Say it:} say \emph{maṭār}
  \item \emph{Recall:} say the Arabic for a corridor, then the Arabic for a public square, then say \emph{maṭār} again
\end{itemize}

\begin{tcolorbox}[breakable,colback=teal!4,colframe=teal!35!black,title={Before you move on}]
What does \ar{مطار} mean? (\textquotedblleft{}An airport\textquotedblright{}.) Say it once more.
\end{tcolorbox}
\section[\texorpdfstring{\ar{ميناء} (mīnāʾ)}{ميناء (mīnāʾ)}]{\ar{ميناء} (mīnāʾ) — a port}

\label{lesson:AR-C100-mina}



\begin{quote}
Before the new one: say the Arabic for a public square, then the Arabic for an airport.
\end{quote}

\subsection*{You'll want to know: \ar{ميناء}}
\textbf{\ar{ميناء}} — \emph{mīnāʾ} — \textquotedblleft{}a port\textquotedblright{}.

\begin{cousinweb}
\textbf{\ar{ميناء}} (\emph{mīnāʾ}) is listed by Lane under \emph{w-n-y}, with the note that it may instead be an ancient Egyptian word; the origin is not settled.
\end{cousinweb}

\begin{grammarlens}[title={What you've built}]
One more word for this chapter.
\end{grammarlens}

\subsection*{Your turn}
\begin{itemize}
\raggedright
  \item \emph{Say it:} \emph{mīnāʾ}
  \item \emph{Say it:} \emph{mīnāʾ}, once more
  \item \emph{Say it:} say \emph{mīnāʾ}
  \item \emph{Recall:} say the Arabic for a public square, then the Arabic for an airport, then say \emph{mīnāʾ} again
\end{itemize}

\begin{tcolorbox}[breakable,colback=teal!4,colframe=teal!35!black,title={Before you move on}]
What does \ar{ميناء} mean? (\textquotedblleft{}A port\textquotedblright{}.) Say it once more.
\end{tcolorbox}
